\section{A signed factorial output: the 508-variable refinement}
\label{sec:508}

The preceding 512-variable equation is a useful control because it changes
only three rows of the final parent. A further change at the arithmetic
source recovers the factorial output's sign from the equations. Its three
private sign witnesses can therefore be removed. Sharing the existing square
also permits one additional exact projection in the rebuilt system.

\subsection{The changed source}
Keep $F$ as an unrestricted integer, retain the cofinal horizon
$t=16+u^2$, and change the arithmetic definitions to
\[
 \rho=F^2t^2,\qquad
 R=2+t+x+y+F^2,\qquad K=2R,\qquad
 \boxed{y=V(V-1).}
 \tag{12}\label{eq:signedchanges}
\]
The power request remains $V=(\rho+1)^t$, but its meaning is established
only after the initial sign argument. The existing source already has the
wire for $F^2$ because its reciprocal product uses
$B=\binom{F^2+t}{t}$. Thus \eqref{eq:signedchanges} adds no fresh square
coordinate. The formula for $y$ is quadratic in the existing $V,y$.

Deleting the factorial guard is not justified by silently assuming $F=t!$.
Indeed, before the power requests are interpreted, $F$ and $V$ may have
either sign. The crucial simplification is that every integer $v$ satisfies
$v(v-1)\ge0$: either $v\ge1$ or $v\le0$. Consequently the literal new
modulus row gives
\[
 y=V(V-1)\ge0,\qquad y+1\ge1
 \tag{13}\label{eq:consecutivepositive}
\]
before any indexed-Pell or power interpretation. This use of discreteness
does not hold for arbitrary real $V$; the domain here is $\Z$.
All independent seed indices are therefore positive, in particular
\[
 t,\ K,\ F^2t^2+t,\ F^2+t,\ y+1,\ x+1\;>0.
 \tag{14}\label{eq:earlyy}
\]
The remaining seeds are the retained positive encoded representatives.
The local constructor bound is also positive on those
representatives. Here $x\ge0$ and every unchanged seed guard are retained.
In particular $R\ge2+t>0$ after \eqref{eq:earlyy}.

The unchanged common budget and guarded frame norm give the same parameter
bounds as in the parent. Every seed inequality implies a positive Pell ordinate. Their shared
auxiliary product is nonzero, so its norm establishes the previously
required auxiliary-coordinate bound. The signed indexed-Pell theorem is
therefore licensed in the same manner as in the parent. The old positivity
certificate is replaced by these actual seed formulas; it is not cited
unchanged after the deletion of its factorial premise. No exceptional
zero-output case or parity restriction on the horizon is needed.

\subsection{Recovering powers and the factorial in the correct order}
Let $D$ denote the proposed factorial binomial output. The retained natural
remainders give $D\ge1$ and a common radix $q\ge2$. The request $q^t$
orients the shared $t$-indexed Pell pair. Thus
\[
 V=(F^2t^2+1)^t\ge1.
\]
The raw ceiling rows give $DF=V+s$, where $s\ge0$. Since $D\ge1$,
this forces $F\ge1$ before any factorial value is assumed. Since $t\ge16$,
$V\ge F^2t^2+1\ge F+2$, and therefore
\[
 y-F-1=V(V-1)-F-1\ge0.
\]
The old endpoint slack can now be recovered.

The reciprocal raw index row is $C+F^2+1=B i_q$, with $C\ge1$ and
$B\ge0$ still guarded. It gives $B>0$. In the signed power request
$N=B^K$, the only possible inverse case is $B=1,N=-1$. But then the
other base at exponent $K$ is $N+V+1\ge2$. That request orients the
$K$-indexed pair and excludes the inverse case. The common radix has its
intended value
\[
 q=N(N+V+1)^K>V^2.
\]
Every remaining seed is similarly oriented by a request at base $q$ or
$q+1$. This establishes all power meanings without first interpreting the
factorial or the binomials whose digits are under validation.

Now $\rho=F^2t^2\ge t$, and for $0\le j\le t$,
\[
 \binom{\rho+t}{j}\le\binom{\rho+t}{t}
 \le(\rho+1)^t=V<q.
\]
The first digit graph therefore really gives $D=\binom{\rho+t}{t}$.
The exact factorial argument is recorded next, rather than inferred from
finite sample calculations.

\begin{lemma}[Self-bounding factorial at the squared radius]
Let $F,t$ be positive integers, $\rho=F^2t^2$,
$V=(\rho+1)^t$, and $D=\binom{\rho+t}{t}$. Then
\[
 F-1<V/D\le F\quad\Longleftrightarrow\quad F=t!.
\]
\end{lemma}
\begin{proof}
Write $G=t!$. The exact identity is
\[
 \frac VD=G\prod_{i=1}^t
   \left(1-\frac{i-1}{\rho+i}\right).
\]
All factors lie in $(0,1]$, and
\[
 \sum_{i=1}^t\frac{i-1}{\rho+i}
 <\frac{1}{2F^2}.
\]
The elementary product bound gives
$G(1-1/(2F^2))<V/D\le G$.
If $F-1<V/D\le F$, then $G>F-1$, so $G\ge F$.
On the other hand,
\[
 G<\frac{F}{1-1/(2F^2)}
   =F+\frac{F}{2F^2-1}\le F+1,
\]
and integrality gives $G\le F$. Thus $F=G$.
Conversely, if $F=G$, then
$F-1<F-1/(2F)<V/D\le F$; when $F=1$ the same strict inequality
also holds. The ceiling remainder $DF-V$ lies in $[0,D-1]$, giving
the required natural ceiling witnesses.
\end{proof}

The source ceiling rows are precisely the inequality in this lemma, so
$F=t!$. The second digit is also safe: for $0\le j\le t$,
\[
 \binom{F^2+t}{j}\le(F^2+t)^t\le(\rho+1)^{2t}=V^2<q.
\]
It recovers the reciprocal product $B$ and $B\ge2$. The later amplified
CRT and derived-label bounds are preserved: after $F\ge1$ the new
$R=2+t+x+y+F^2$ is at least the old $2+t+x+y+F$. The same local
majorant exponent, factorial growth floor, representative bounds, and
strengthened CRT precision therefore apply. No preceding bootstrap step
used that later factorial growth estimate.

\subsection{Completeness and the literal polynomial}
Choose the cofinal padded horizon as before and set $F=t!$,
$\rho=F^2t^2$, and $V=(\rho+1)^t$. A slightly stronger pending-product
bound is useful. Put $a=Ft\ge16$, $L=a+1$, and $B_0=a^2+1$.
There are at most $t+1$ pending weights, each at most $L$, so
$S_i\le L^{t+1}$. Since $a\ge3$ and $t\ge2$,
\[
 B_0>(a-1)L,\qquad (a-1)^t\ge(a-1)^2\ge L.
\]
It follows, using integer powers only, that
\[
 V=B_0^t>L^t(a-1)^t\ge L^{t+1}\ge S_i.
\]
Now $V\ge3$, so the prescribed modulus $y=V(V-1)$ is at least $V$
and exceeds every accumulator. Also $F\mid(V-1)$ because
$V\equiv1\pmod F$, and hence $F\mid y$. We have already established
$y\ge F+1$. These are exactly the inherited CRT requirements; the earlier
sufficient convenience $y>V^2$ is unnecessary. The first CRT is valid. The larger new
$R$ satisfies every subsequent lower bound. The factorial lemma supplies
the ceiling witnesses, the unchanged digit arguments supply their quotients,
and simultaneous Pell completeness supplies the common frame. Thus every
finite refutation gives a zero of the new arithmetic system.

The source is then factored, projected through monic integral definitions,
and translated by the same recorded pipeline as its matched parent. The
principal result is
\[
 \boxed{P_{508}(x_0,\ldots,x_{507})
       =\sum_{j=0}^{307}Q_j(x_0,\ldots,x_{507})^2=0.}
 \tag{15}\label{eq:P508}
\]
The full $Q_j$ are printed in Appendix~\ref{app:principal-residuals} and supplied in
\file{direct/principal/equation.json}. Their full collected expansion is
\file{direct/principal/expanded_polynomial.json}.

\begin{theorem}[Principal unrestricted-integer equation]
The polynomial in \eqref{eq:P508} has 508 occurring integer variables,
308 quadratic residuals, exact degree four, 6,545 collected monomials, and
coefficient height 996,521. ZFC proves
\[
 \bigl[\forall\mathbf x\in\Z^{508}\;P_{508}(\mathbf x)\ne0\bigr]
 \longleftrightarrow\Con(\ZFC).
\]
\end{theorem}
\begin{proof}
The product of consecutive integers establishes the modulus sign, and hence
all seed signs, without assuming a factorial or power meaning. The ordered recovery gives the
same bounded local proof-certificate semantics, whose soundness gives
\eqref{eq:calibration}. The constructive paragraph above proves the
converse from any finite refutation. Every source modification is
identified by an exact quadratic-row certificate. The subsequent monic
eliminations and integer translations have recorded integral inverse maps,
so they preserve integer-zero existence. Finally a sum of integral squares
vanishes exactly when all its residuals vanish. Exact coefficient expansion
establishes the finite metrics, including a nonzero degree-four part.
All quantified mathematical steps are formalizable in ZFC, yielding
\eqref{eq:target} for $D=\Z$ and $n=508$.
\end{proof}

The matched parent has 517 variables, 310 rows, 6,774 terms, and coefficient
height 3,904,267. This refinement improves all four of those quantities
while retaining degree four and unrestricted integer unknowns. It removes
about $1.74\%$ of the variables and $3.38\%$ of the monomials; coefficient
height falls by about $74.48\%$. Those are measured improvements in this
particular direct construction, not a claim that 508 is the smallest
possible integer-variable count.
